\clearpage
\fancyhead[L]{\small JUNO. DISCOVERY.}
\hypertarget{native-258-lower}{}
\begin{center}
SOME NEWS ABOUT THE NEW PLANET.

\rule{0.75\linewidth}{0.4pt}

Braunschweigisches Magazin. No. 40. 1804, 6 October.

\rule{0.75\linewidth}{0.4pt}
\end{center}

The recent discovery of a new moving star at \emph{Lilienthal}, which has already been reported in the newspapers, has aroused so much public curiosity that people who take an interest in the progress of science will probably welcome some more definite news of this remarkable discovery. This is especially true because, now that
\SourceLeaf{259}
the star’s course has been observed for three weeks, fairly reliable conclusions can already be drawn about its true orbit. Most readers of these pages would otherwise see those results only much later.

Mr. \textsc{Harding}, an associate of the renowned Councillor of Justice \textsc{Schröter} in Lilienthal, had already been working for half a year on an astronomical atlas of a special kind. It is intended to cover only those parts of the fixed-star sky in which the two newly discovered planets \emph{Ceres} and \emph{Pallas} can be seen. Finding and observing these small celestial bodies requires very detailed star charts that include every star down to at least the eighth magnitude. All the existing charts of the heavens, without exception, fall far short of what is needed. Mr. \textsc{Harding}’s undertaking is therefore well deserving of recognition. Some excellent detailed charts that he has already contributed to Mr. von \textsc{Zach}’s Monatliche Correspondenz attest to the great care and tireless diligence with which he is carrying it out. Mr. \textsc{Harding} relies chiefly on \textsc{Lalande}’s Histoire céleste française, an invaluable work containing observations made over 10 years at the observatory of the military school in Paris, chiefly by \textsc{Lalande}’s nephew, \textsc{Michel Lefrançais Lalande}, and covering 50000 fixed stars. Mr. \textsc{Harding} records in his charts every one of those stars that falls within the zones he is surveying. But he first looks for each star himself in the sky, to make sure that it really is there and to correct any errors of observation, writing, or printing that are unavoidable in so enormous a work. Nor does Mr. \textsc{Harding} stop there. Wherever he encounters fainter fixed stars, down to at least the eighth magnitude, that are absent from the Histoire céleste, he determines their positions himself and adds them to his charts. A great many such stars remain.

It was during just such a survey of a region of the sky that Mr. \textsc{Harding} came across a star of the eighth magnitude in Pisces on 1 September. It was absent from the Histoire céleste. Since that happens countless times, he suspected nothing unusual and entered it in his charts as a fixed star, provisionally judging its position by eye alone. But when he looked again three days later, there was no longer a star in that spot. A very similar one appeared nearby, about half a degree to the southwest, where he could not remember seeing a star on 1 September. This aroused strong suspicion, though it still warranted no definite conclusion, since it might merely have been due to a change in brightness, quite common among faint fixed stars, or to an error. Finally, on 5 September both former spots were empty, and a spot a proportionate distance to the southwest was occupied. Mr. \textsc{Harding} could now have no further doubt that the stars seen on the three evenings were one and the same. He therefore made a formal observation of the moving star he had discovered and continued his observations on the following evenings. They confirmed a perfectly regular retrograde motion that also carried the star southward. He informed astronomers elsewhere of his important discovery.

I received the news on the 11th, when the sky here was completely overcast in the evening. On the 12th, however, it was very clear. I found \textsc{Harding}’s star at once and determined its position by a good observation. From then until today (30 Sept.), I have been able to observe it on eleven more evenings, most recently on the 28th. Its motion is still retrograde and southward. The arc it has traced so far is so close to a straight line that the curvature is scarcely noticeable, but its speed has increased perceptibly. These phenomena agree very well with the supposition that this celestial body really does describe a planetary orbit. Taken in these general terms, however, they naturally warrant only conjectures at this stage. Only a
\SourceLeaf{260}
calculation carried out with as much precision and caution as so delicate and tricky an investigation demands can give us certainty.

In appearance this star is very much like Ceres and Pallas, though it is slightly brighter than they are. I estimate that it is about as bright as a fixed star of the $7^{\mathrm{th}}$ to $8^{\mathrm{th}}$ magnitude. Even on 18 September, when the nearly full Moon was close by, I was able to see and observe it with a \textsc{Short} reflecting telescope without the slightest effort. This star shows no trace at all of the haze seen around comets, even through the large instruments at Lilienthal.

Of the other astronomers elsewhere who are already observing this star, I know so far only of Dr. \textsc{Olbers}, Mr. von \textsc{Zach}, and Professor \textsc{Bode}. The first began his observations on the 7th, Mr. von \textsc{Zach} on the 13th, and Mr. \textsc{Bode} on the 21st. Since the star will remain visible until at least January 1805, I confidently hope that it will still be possible to observe it at all European observatories. Although this remarkable celestial body has been observed for only a short time, I have already attempted to make a first approximate determination of the principal dimensions of its true orbit, mostly from my own observations. This serves in part to satisfy, for the time being, the very natural curiosity of the public. But such work has a more substantial use in that it enables us to predict the star’s course for some time \emph{ahead}. This lets us give more distant astronomers, who may receive reports of the course already travelled too late to find the star directly from them, a rough indication of where they should look and perhaps spare them a fruitless search.

The remarkable result of this investigation was that \textsc{Harding}’s star, like Ceres and Pallas, describes a planetary orbit between Mars and Jupiter. Its period of revolution and mean distance from the Sun are almost (or exactly) the same as theirs. Its eccentricity is slightly greater even than that of Pallas, whereas its orbital inclination is much smaller than that of Pallas and only slightly greater than that of Ceres. The results were a period of revolution of 1662 days, a mean distance of 2.746, an eccentricity of 0.287, an orbital inclination of 13 degrees 35 minutes, a longitude of aphelion of 239 degrees 14 min., a longitude of the ascending node of 171 degrees 16 min. The mean longitude in the orbit on 5 Sept. at noon on the meridian of the Seeberg Observatory was 20 degrees 39 min.

On the strength of these results, it can already be regarded as settled that \textsc{Harding}’s star, just like the other planets, will have its regular periods of visibility and invisibility and will henceforth be a constant object of our observations. It therefore fully deserves the rank of a planet. To put it on the same footing as the older planets, Mr. \textsc{Harding} has already proposed the name \emph{Juno} for it. There is no doubt that all astronomers will respect the discoverer’s right and accept his choice of name. Juno’s sceptre bearing a star has been proposed as her symbol.

Juno now lies within the easily recognisable quadrilateral formed by the four stars in Pisces numbered 27, 29, 30, 33 in \textsc{Flamsteed}’s catalogue. In a clear sky a person with good eyesight can see these stars without a telescope. Juno will continue moving in its present direction for some time, though its speed is already decreasing again. In the latter half of October the path will curve ever more noticeably to the southeast. Around 30 October the retrograde motion becomes direct, and around 17 November, when Juno will have reached a southern declination of $11\frac{1}{4}$ degrees, that declination begins to decrease again and the path curves back northward.

\SourceLeaf{261}
There may be a reader here and there with instruments and skill enough to try to find this new planet and follow it along its course. I therefore add its exact past and future positions, the former at weekly intervals only, the latter every 3 days. All are for midnight. [This ephemeris is an extract from the one printed above on p. 253.]
